\subsection*{The continuous scattering picture}
The starting point is the vector-Maxwell volume-integral equation for an isotropic, nonmagnetic dielectric. For one illumination, write the induced current, or equivalently a consistently scaled contrast source, schematically as
\[
j(\boldsymbol r)=X(\boldsymbol r)\left[E^{\rm inc}(\boldsymbol r)
+\int_D G_D(\boldsymbol r,\boldsymbol r')j(\boldsymbol r')\,d\boldsymbol r'\right].
\]
Here $G_D$ maps source currents to the internal electric field, $X$ is the local material response, and the bracket is the total field $E^{\rm tot}$. The singular self-interaction is understood in the chosen Maxwell volume-integral convention; the implemented cell self-response is specified below. The physical constants associated with the source convention are included consistently in $X$ and the Green maps.

A small admissible change $\delta\chi$ changes both the local material response and the field produced by the currents. Differentiating at the current material state gives
\[
\delta j(\boldsymbol r)=DX[\delta\chi](\boldsymbol r)E^{\rm tot}(\boldsymbol r)
+X(\boldsymbol r)\int_DG_D(\boldsymbol r,\boldsymbol r')\delta j(\boldsymbol r')\,d\boldsymbol r'.
\]
The first term injects the material perturbation into the current equation. The second rescatters the resulting current perturbation. Moving rescattering to the left defines
\[
L=I-XG_D,\qquad \mathcal B\delta\chi=DX[\delta\chi]E^{\rm tot},
\qquad L\delta j=\mathcal B\delta\chi.
\]
For receiver channel $m$, including its polarization readout, the observation equation is
\[
\delta y_m=\int_DG_{S,m}(\boldsymbol r_m,\boldsymbol r')\delta j(\boldsymbol r')\,d\boldsymbol r'
=:(S\delta j)_m.
\]
At a well-posed material state where $L$ is invertible, $\delta y=SL^{-1}\mathcal B\delta\chi$. In physical material coordinates $\delta\chi=Q_Ts$, define $B=\mathcal BQ_T$; the material Jacobian is then $J=SL^{-1}B$. Receiver normalization or noise weighting, when used, is included in $S$ and the data inner product.

Carrying the spatial integrals through each elimination, adjoint pullback, and multi-illumination calculation would obscure this structure. We therefore use operator notation from this point onward. Three levels remain distinct: the continuous contrast-source equation supplies the physical interpretation; $L\delta j=Bs$, $\delta y=S\delta j$ supplies the abstract derivative structure; and the implemented DDA realizes the local response through the contrast-dependent cell polarizability $a(\chi)$. Its $B$ contains $a'(\chi)$ and the three-component total field, as shown in (S2)--(S3).

\subsection*{Why an adjoint returns information to material space}
The adjoint is defined by the data and physical material inner products:
\[
\langle Js,q\rangle_{\rm data}=\langle s,J^*q\rangle_{\rm material},
\qquad J^*=B^*L^{-*}S^*.
\]
Here $^*$ denotes adjoints in the realified data, current, and physical material spaces. Equivalently, in metric-normalized complex arrays with real material coordinates, $J^*q=\operatorname{Re}(B^HL^{-H}S^Hq)$. The coordinate convention below makes this real-part pullback explicit.

Whereas $J$ predicts the measurement change caused by a material direction, $J^*$ maps a receiver residual or test vector to material directions that could explain it. Read the adjoint product from right to left: $S^*q$ puts receiver information into the current equation; $L^{-*}$ propagates it through the adjoint multiple-scattering system; $B^*$ pulls that response back through the local-field and material factors. This is an adjoint source calculation with the stated metrics, rather than a new physical illumination experiment.

This return path prepares the two material factors used below. $J_0^*Y_C$ pulls the added observation through the original material sensitivity, while $N_C^*$ returns information through the added material-injection path. Both enter when the material normal equation changes.

\subsection*{How a retained current basis produces feedback}
Approximate a current perturbation by $\delta j(\boldsymbol r)=\sum_i v_i(\boldsymbol r)\alpha_i$. Galerkin testing requires the current-equation residual to be orthogonal to the retained basis $V$. Hence $V^*LV\alpha=V^*Bs$, and, for an invertible retained block,
\[
R_0=V(V^*LV)^{-1}V^*,\qquad \delta j_0=R_0Bs.
\]
$R_0$ is the reduced current solution operator. Add orthonormal directions $C$ outside $V$. Eliminating the old coefficients leaves a Schur equation for the new ones. Back-substitution produces
\[
\bar C=(I-R_0L)C=C+R_0XG_DC,\qquad
Y_C=S\bar C,\quad N_C=C^*(B-LR_0B),\quad\Gamma_C=C^*L\bar C.
\]
The added current directly radiates through $SC$ and also excites retained currents, whose response returns through $SR_0XG_DC$. Their sum $S\bar C$ is the dressed observation. Thus $\bar C$ follows from solving the retained feedback loop, with no empirical correction factor. Proposition 1 in the main manuscript gives the elimination and invertibility conditions. The next subsection fixes the coordinates and the exact DDA realization needed to evaluate these operators.

\subsection*{Physical coordinates and the implemented vector-Maxwell model}

All material normal equations are over real Hilbert spaces, with material coordinates \(s\in\mathbb R^d\). For complex material contrast, the real and imaginary grid values are stacked, and the real material expansion satisfies \(Q_T^\top M_\chi Q_T=I\). A complex-output map \(F\) of real material coordinates has real adjoint \(v\mapsto\operatorname{Re}(F^Hv)\); equivalently, use its fully realified matrix. For a complex current operator,

\[
\mathbb R(F)=\begin{bmatrix}\operatorname{Re}F&-\operatorname{Im}F\\
\operatorname{Im}F&\operatorname{Re}F\end{bmatrix}.
\]

For \(P\) illuminations, \(L\), \(S\), \(V\), and \(C\) lift blockwise, while the \(B_p\) stack vertically and share the same \(s\). In particular, \(C_{\rm blk}=I_P\otimes\mathbb R(C)\) has \(2Pr_c\) columns. The material curvature is the sum $J^TJ=\sum_p\operatorname{Re}(J_p^HJ_p)$ over shared material coordinates; no cross-illumination term is needed in this material normal matrix. In contrast, the output covariance $JJ^T$ has off-diagonal illumination blocks. It cannot be replaced by unrelated per-illumination inverse problems.

The current and material metrics are normalized before all orthogonalizations. Orthogonal rotations of an admissible physical tangent preserve the assertions. Different Gaussian and voxel spans are different admissible tangents. Their numerical comparison does not establish an invariant ranking across different physical spaces.

The actual reused DDA models an isotropic nonmagnetic dielectric, \(\chi=\epsilon_r-1\), on cell centers in a box of side 1.5. A cell has volume \(v=(1.5/n)^3\). With \(\boldsymbol n=\boldsymbol r/\|\boldsymbol r\|\), the off-diagonal electric dyadic is

\[
G(\boldsymbol r)=\frac{e^{ikr}}{4\pi r}
\left[k^2(I-\boldsymbol n\boldsymbol n^\top)
+\left(\frac{ik}{r}-\frac1{r^2}\right)(I-3\boldsymbol n\boldsymbol n^\top)\right]. \tag{S1}
\]

Self blocks are excluded from \(G_{\rm off}\). The implemented cell polarizability and derivative are

\[
a(\chi)=\frac{3v\chi}{\chi+3-i\,3c_kv\chi},\qquad
c_k=\frac{k^3}{6\pi},\qquad
a'(\chi)=\frac{9v}{(\chi+3-i\,3c_kv\chi)^2}. \tag{S2}
\]

For dipole moments \(p_p\), the normalized current is \(c_p=p_p/\sqrt v\). The state, total field, tangent, and measurement equations are

\[
\begin{aligned}
L&=I-D_aG_{\rm off},&
Lc_p&=D_ae_p^{\rm inc}/\sqrt v,\\
E_p&=e_p^{\rm inc}+G_{\rm off}(\sqrt v\,c_p),&
L\delta c_p&=D_{a'E_p}Q_Ts/\sqrt v,\\
y_p&=\sqrt v\,G_{\rm rec}c_p.
\end{aligned} \tag{S3}
\]

The material-to-vector map in the tangent multiplies each cell's scalar material perturbation by its three field components. It is not a scalar diagonal on an unrelated vector space. Polarizability is not contrast; dropping \(a'(\chi)\) changes the Jacobian. At zero contrast \(a'=v\), so zero contrast does not remove the material tangent. Equation (S2) incorporates Clausius--Mossotti and radiation-reaction terms in one expression; the final denominator, not an intermediate rearrangement at \(\chi=-3\), specifies its singularities. These formulas document the reused code. Reference~\cite{r11} supplies DDA background, not an audited derivation of this exact implementation.

There are six transverse plane-wave illuminations and 64 receivers on a sphere of radius 5, with two transverse readout functionals each. The receiver map retains the dyadic near-field terms. Measured scattered fields and tangent actions share the scalar normalization \(\|y_{\rm clean}\|\); no estimated noise-covariance whitening was used in the experiments. This clean norm remains known and fixed in the synthetic noise checks.

The material metric is \(vI\) for each complex component. Gaussian columns are \(\exp[-\|x_i-\xi_\nu\|^2/(2\times0.28^2)]\), with \(\xi_\nu\in\{-0.45,0,0.45\}^3\). An SVD of the volume-weighted columns produces the orthonormal physical tangent. Voxel expansion is \(\delta\chi=(s_{\rm re}+is_{\rm im})/\sqrt v\). Smooth Gaussian profiles define the internal truth, with object-dependent weights and a constant background; their fixed-box support is truncated at the outer boundary. This paragraph describes the smooth-object subset; the sharp-interface protocols are retained in the historical supplement archive.

